% Generated draft: requires mathematical and source review.
\documentclass[11pt]{article}
\usepackage[margin=25mm]{geometry}
\usepackage{fontspec}
\setmainfont{Latin Modern Roman}
\usepackage{amsmath,amssymb,mathtools,unicode-math}
\setmathfont{Latin Modern Math}
\usepackage{enumitem,needspace,xurl,hyperref}
\setlength{\parindent}{0pt}
\setlength{\parskip}{5pt}
\setlength{\emergencystretch}{3em}
\newcommand{\N}{\mathbb{N}}
\newcommand{\Z}{\mathbb{Z}}
\newcommand{\Q}{\mathbb{Q}}
\newcommand{\R}{\mathbb{R}}
\newcommand{\C}{\mathbb{C}}
\newcommand{\F}{\mathbb{F}}
\newcommand{\A}{\mathbb{A}}
\newcommand{\PP}{\mathbb{P}}
\newcommand{\E}{\mathbb{E}}
\newcommand{\eps}{\varepsilon}
\newcommand{\dd}{\,\mathrm{d}}
\begin{document}
\section*{UnsolvedMath ID: 700193}
\textbf{Source problem code:} catalogue entry
\subsection{Definitions and mathematical statement}
For a group \(H\), define its derived series by \(H^{(0)}=H\) and
\[
H^{(n+1)}=\bigl\langle [x,y]:x,y\in H^{(n)}\bigr\rangle,
\qquad [x,y]=x^{-1}y^{-1}xy.
\]
The group \(H\) is \emph{soluble} if \(H^{(d)}=\{1\}\) for some integer \(d\geq 0\). A group is \emph{locally soluble} if every subgroup generated by finitely many elements is soluble.

A subgroup \(H\leq G\) is \emph{ascendant in \(G\)} if there are an ordinal \(\alpha\) and an increasing chain of subgroups \((H_\beta)_{\beta\leq\alpha}\) such that
\[
H_0=H,\qquad H_\alpha=G,\qquad
H_\beta\trianglelefteq H_{\beta+1}\quad(\beta<\alpha),
\]
and
\[
H_\lambda=\bigcup_{\beta<\lambda}H_\beta
\quad\text{for every nonzero limit ordinal }\lambda\leq\alpha.
\]

The question is whether the following implication holds for every group \(G\): if there is a family \((H_i)_{i\in I}\) of soluble subgroups, each ascendant in \(G\), with
\[
G=\left\langle\bigcup_{i\in I}H_i\right\rangle,
\]
then
\[
\forall\,\text{finite }F\subseteq G\quad
\exists\,d\in\mathbb Z_{\geq0}\quad
\langle F\rangle^{(d)}=\{1\}.
\]
There is no finiteness assumption on \(G\) or on the generating family, and no uniform bound on the derived lengths of the \(H_i\) or on the lengths of their ascendant chains is assumed.
\subsection{Short English statement}
B. I. Plotkin asks whether every group generated by soluble subgroups that are ascendant in the whole group must be locally soluble, meaning that each finitely generated subgroup is soluble. The supplied record labels the problem open.
\subsection{Sources}
\begin{itemize}
\item Local UnsolvedMath registry, permanent ID \(700193\), source problem code \texttt{catalogue entry}, titled \emph{Kourovka Notebook Problem 3.44}. The supplied record labels its status as open.
\item \emph{The Kourovka Notebook (Unsolved Problems in Group Theory)}, original problem 3.44, third issue (1969); proposed by B. I. Plotkin.
\item The supplied record identifies the twenty-first issue edition (2026), \texttt{arXiv:1401.0300v48}. Its collection description names E. I. Khukhro and V. D. Mazurov as editors. These bibliographic details are supplied attribution and have not been independently verified here.
\end{itemize}
\end{document}
